\section*{Bab \ref{ch:g4:symmetry} --- Simetri}

\begin{solution}{exo:g4:symmetry:1}
Persegi: ya ($4$ sumbu). Persegi panjang: ya ($2$). Segitiga
sembarang: tidak ada lipatan yang berhasil. Lingkaran: ya --- setiap garis yang melalui pusatnya.
\end{solution}

\begin{solution}{exo:g4:symmetry:2}
Persegi: $4$ (dua garis tengah, dua diagonal). Persegi panjang: $2$
(hanya garis tengahnya). Segitiga sama sisi: $3$ (satu melalui tiap
\omterm{def:g4:geometry:polygon}{titik sudut}).
\end{solution}

\begin{solution}{exo:g4:symmetry:3}
Ketika dilipat sepanjang diagonalnya, kedua paruh segitiga itu punya
bentuk yang sama tetapi tidak saling menutupi --- yang satu
``terbalik arah'': jiplakannya mendarat di samping gambar, bukan di
atasnya. Jadi diagonalnya bukan sumbu.
\end{solution}

\begin{solution}{exo:g4:symmetry:4}
Huruf L bayangannya berada $2$ petak di sisi \emph{yang lain} dari
sumbu, tertulis terbalik (huruf L cermin), dengan ukuran yang sama.
\end{solution}

\begin{solution}{exo:g4:symmetry:5}
Gambar yang lengkap menunjukkan bendera itu beserta pantulannya yang
terbalik di bawah sumbu, tiang di bawah tiang.
\end{solution}

\begin{solution}{exo:g4:symmetry:6}
$P$ tidak berpindah: melipat sepanjang sumbu membiarkan setiap titik
\emph{pada} sumbu tetap di tempatnya --- $P$ adalah bayangannya sendiri.
\end{solution}

\begin{solution}{exo:g4:symmetry:7}
Bersumbu tegak: A, H, O, T. Bersumbu mendatar: B, E, H, O.
Keduanya: H, O. Bukan keduanya: F, N.
\end{solution}

\begin{solution}{exo:g4:symmetry:8}
Ya: segitiganya sama kaki, dan sumbunya melalui ujung di antara
kedua sisi $5$ cm dan tengah alas $3$ cm. Lipatan itu
mempertukarkan kedua sisi yang sama panjang.
\end{solution}

\begin{solution}{exo:g4:symmetry:9}
Setiap bangun dengan tepat dua sumbu berhasil: persegi panjang (yang
bukan persegi), tanda tambah dengan lengan dua panjang yang berbeda
(mendatar panjang, tegak pendek), atau bangun lonjong. Kedua
sumbunya selalu berpotongan di pusat bangun itu.
\end{solution}

\begin{solution}{exo:g4:symmetry:10}
Rumah yang lengkap adalah paruh kiri beserta kembaran cerminnya:
dinding selebar dua kali separuh dinding, di bawah atap segitiga
yang utuh. Setiap panjang di kanan sama dengan pasangannya di kiri
--- simetri menyalin panjang dengan tepat.
\end{solution}

\begin{solution}{exo:g4:symmetry:11}
Gambarnya mendukung kedua arah: segitiga sama kaki terlipat
sepanjang garis yang melalui ujungnya (dua sisi sama panjang $\to$
sumbu); segitiga sama sisi terlipat dengan tiga cara; segitiga
sembarang tidak punya sumbu (tidak ada sisi yang sama panjang $\to$
tidak ada sumbu). Pernyataan ganda Aya benar --- itu akan dibuktikan
pada \cref{ch:g6:symmetry}.

\end{solution}
